% Estas slides tienen que abrirse con el programa pdfpc que soporta videos embebidos
% el comando es: pdfpc -g slides.pdf
% para los videos se requiere ubuntu-restricted-extras
% para la bibliografía se requiere biber y configurar texstudio


%\documentclass[compress,handout]{beamer}
\documentclass[compress]{beamer}


% Theme customization
\setbeamertemplate{itemize item}[rectangle] % configure itemize
\setbeamertemplate{itemize subitem}[circle] % configure itemize
\setbeamertemplate{itemize subsubitem}[triangle] % configure itemize
\setbeamertemplate{navigation symbols}{} % remover simbolos de navegacion de las slides
\usefonttheme[onlymath]{serif} % simbolos matematicos en serif (Como es en latex original)

\setbeamertemplate{blocks}[rounded] % blocks corners rounded
\setbeamercolor{block body}{bg=blue!12,fg=black} % color of blocks

% Latex packages
\usepackage{pdfpc-commands} % pdfpc movie commands. Requires to be downloaded manually from https://github.com/pdfpc/pdfpc/issues/328#issuecomment-437397536
\usepackage[utf8]{inputenc}
\usepackage[T1]{fontenc} % para copiar acentos en español del pdf y permite acentos en las notas
\usepackage[spanish]{babel}
\usepackage[binary-units=true,per-mode = symbol]{siunitx} % para manejar las unidades
\usepackage{multirow}
\usepackage{graphicx}
\usepackage{xcolor}
\usepackage{amsmath} % bmatrix
\usepackage[caption=false]{subfig} % caption = false elimina la palabra "Figura" del caption
\setbeamertemplate{caption}{\raggedright\insertcaption\par} % elimina la palabra "Figura" del caption
\usepackage{import} % para el comando import (se usa para pdf_tex)
\captionsetup[subfigure]{labelformat=empty} % remover el indice del caption de la subfigura
\usepackage{booktabs} % \toprule \midrule \bottomrule
\usepackage[overridenote]{pdfpc} % requires to download manually pdfpc.sty package from https://www.ctan.org/pkg/pdfpc
\usepackage[backend=biber]{biblatex} % set biber to format references. Must configure Biber in Texstudio
\usepackage{csquotes} % to remove warning triggered by biblatex and babel

% Color commands for annotations
\newcommand\TODO[1]{\textbf{\textcolor{red}{#1}}} %  TODO notes

% add bibliography resource
\renewcommand*{\bibfont}{\footnotesize} % change bibliograhy size
\bibliography{./src/bibliography.bib}

% Graphic paths
\graphicspath{{./images/}}

% add math preamble
%\usepackage{amsmath}
\usepackage{amssymb}
\usepackage{amsopn}
\usepackage{mathtools}

% math
\renewcommand{\vec}[1]{\boldsymbol{\mathbf{#1}}}
\newcommand{\norm}[1]{\lVert#1\rVert}

% Declare arg max and arg min functionss
\DeclareMathOperator*{\argmax}{arg\,max}
\DeclareMathOperator*{\argmin}{arg\,min}

% Homogeneous decoration function
\newcommand{\homo}[1]{\dot{#1}}


% Declare projection as math function
\DeclareMathOperator{\proj}{proj}
\newcommand{\fromCoord}[2]{{#1}_\mathrm{#2}}
\newcommand{\toCoord}[2]{\prescript{\mathrm{#2}}{}{#1}}
\newcommand{\worldCoordSystem}{\mathrm{w}}
\newcommand{\bodyCoordSystem}{\mathrm{B}}
\newcommand{\cameraCoordSystem}{\mathrm{c}}
\newcommand{\point}{\vec{p}}
\newcommand{\worldPoint}{\toCoord{\point}{\worldCoordSystem}}
\newcommand{\imagePoint}{\vec{u}}
\newcommand{\cameraPoint}{\toCoord{\point}{\cameraCoordSystem}}
\newcommand{\homoWorldPoint}{\toCoord{\homo{\point}}{\worldCoordSystem}}
\newcommand{\homoImagePoint}{\homo{\imagePoint}}
\newcommand{\homoCameraPoint}{\toCoord{\homo{\point}}{\cameraCoordSystem}}
\newcommand{\measurement}{\vec{z}}
\newcommand{\prediction}{\hat{\vec{z}}}
\newcommand{\seMatrix}{\vec{\xi}}
\newcommand{\transform}[2]{\toCoord{\fromCoord{\seMatrix}{#2}}{#1}}
\newcommand{\pointCoord}[1]{\toCoord{\point}{#1}}
\newcommand{\rotation}{\vec{R}}
\newcommand{\rotationCoord}[2]{\toCoord{\fromCoord{\rotation}{#2}}{#1}}
\newcommand{\translation}{\vec{t}}
\newcommand{\translationCoord}[2]{\toCoord{\fromCoord{\translation}{#2}}{#1}}
\newcommand{\intrinsicMatrix}{\vec{K}}
\newcommand{\principalPoint}{\vec{c}}
\newcommand{\reprojectionError}{u}
\newcommand{\projectionMatrix}{\vec{P}}
\newcommand{\cameraCenter}{\vec{o}}
\newcommand{\essentialMatrix}{\vec{E}}
\newcommand{\inverse}[1]{{#1}^{-1}}

% Motion model
\newcommand{\position}{\vec{p}}
\newcommand{\orientationQuaternion}{\vec{q}}
\newcommand{\predictedPosition}{\hat{\vec{p}}}
\newcommand{\predictedOrientationQuaternion}{\hat{\vec{q}}}
\newcommand{\linearVelocity}{\vec{v}}
\newcommand{\angularVelocity}{\vec{\omega}}

\DeclareMathOperator{\slerpOp}{slerp}
\newcommand{\slerp}[1]{\slerpOp{\left( #1 \right)}}

% Map structure
\newcommand{\map}{M}
\newcommand{\keyframesSet}{K}
\newcommand{\mapPointsSet}{P}
\newcommand{\observedMapPoints}{O}
\newcommand{\covisibilityKeyframes}{CK}
\newcommand{\localMap}{local\_map}



% Bundle Adjutment
\newcommand{\update}{\vec{\delta}}
\newcommand{\incremental}{\hat{\update}}


% Loop Closure names

% scaled operators and letters to fancy view
\newcommand{\sminus}{\scalebox{0.5}[1.0]{$-$}}
\newcommand{\splus}{\scalebox{0.6}[0.6]{$+$}}
\newcommand{\curr}{c}
\newcommand{\sind}[1]{\scalebox{0.6}[0.6]{$#1$}}
\newcommand{\ind}[1]{\scalebox{0.7}[0.7]{$#1$}}

\newcommand{\keyframe}{\vec{K}}
\newcommand{\bowVector}{\vec{v}}
\newcommand{\lcError}{\vec{\Omega}}
\newcommand{\relativeTransformation}{\seMatrix}
\DeclareMathOperator{\interpolate}{interpolate}

\newcommand{\relativeMotion}{\vec{\delta}}
\newcommand{\groundTruth}[1]{{#1}^{*}}



% definición del operador rot()
\DeclareMathOperator{\rotationOp}{rot}
\newcommand{\getRotation}[1]{\rotationOp{\left( #1 \right)}}

\DeclareMathOperator{\translationOp}{trans}
\newcommand{\getTranslation}[1]{\translationOp{\left( #1 \right)}}









\title{ROS: Robot Operative System}
\author{}
\institute{Universidad Nacional de Rosario}
%\date{\scriptsize{Julio 1, 2021}}
\date{}

\begin{document}

% add title page
\frame{\titlepage}

\section{ROS}
\begin{frame}
	\frametitle{Motivación}
	
	Don’t reinvent the wheel. Create something new and do it faster and better by building on ROS!
	
	\note{https://www.ros.org/blog/why-ros/}
	
	\note{http://wiki.ros.org/es/ROS/Tutoriales}
	
	\note{https://www.youtube.com/playlist?list=PL8dDSKArO2-m7hAjOgqL5uV75aZW6cqE5}
	
\end{frame}

\begin{frame}
	\frametitle{¿Qué es ROS?}
	
	ROS (Robot Operating System) es un kit de desarrollo de software de código abierto para aplicaciones de robótica. ROS ofrece una plataforma de software estándar para desarrolladores de todas las industrias que los llevará desde la investigación y la creación de prototipos hasta la implementación y la producción.
	
	\begin{itemize}
		\item Comunidad global
		\item Utilizado en cursos de robótica, investigación e industria
		\item Acorta los tiempos de producción
		\item Multi-dominio: indoor / outdoor, hogareño / industrial, bajo el agua / espacio
		\item Multi-plataforma: Linux, Windows y macOS.
		\item Open-source
		\item Licencia permisiva (Apache 2.0)
		\item Soporte desde la industria
	\end{itemize}
	
	
	\note{https://www.ros.org/blog/why-ros/}
	
	
\end{frame}

\begin{frame}
	\frametitle{¿Qué es ROS?}
	  Trabajaremos con la versión Ubuntu LTS más actual, esta siempre viene con una versión de ROS estable.
	
	\begin{figure}[!h]
		\centering
		\subfloat[ROS1]
		{
			\includegraphics[width=0.4\columnwidth]{images/ros_version_noetic.png}
		}
		\subfloat[ROS2]
		{
			\includegraphics[width=0.4\columnwidth]{images/ros_version_humble.png}
		}
	\end{figure}

\end{frame}

\begin{frame}
	\frametitle{Configuración de entorno de ROS}
    
    https://docs.ros.org/en/<distro>/Tutorials.html
    
    Instalar ROS
    
    sudo apt install ros-<distro>-<package>
    
    Sourcear los archivos de setup
    
    source /opt/ros/<distro>/setup.bash
    
    Agregar el source al .bashrc
    
    echo "source /opt/ros/<distro>/setup.bash" >> ~/.bashrc
    
    Chequear las variables de entorno
    
    printenv | grep -i ROS
    
    Output:
    
    ROS\_VERSION=2\\
    ROS\_PYTHON\_VERSION=3\\
    ROS\_DISTRO=humble
    
    

    setear el DOMAIN ID:
    
    Los nodos ROS 2 en el mismo dominio pueden descubrir y enviarse mensajes libremente, mientras que los nodos ROS 2 en diferentes dominios no pueden. Todos los nodos ROS 2 utilizan el ID de dominio 0 de forma predeterminada.
    
    echo "export ROS\_DOMAIN\_ID=<your\_domain\_id>" >> ~/.bashrc
	
	
\end{frame}

\begin{frame}
	\frametitle{Creando un paquete en ROS}
	
     Building system: ament
     luego para compilar un paquete: colcon build
     
     
     crear un paquete:
     
     ros2 pkg create test\_pkg --build-type ament\_cmake
	
\end{frame}

\begin{frame}
	\frametitle{Cómo escribir un Publicador y un subscriptor}

\end{frame}


\begin{frame}
	\frametitle{Install turtlebot en ROS}
	
\end{frame}

\begin{frame}
	\frametitle{Publicar velocidades a Turtlebot}
	
\end{frame}

\begin{frame}
	\frametitle{TF}
	Hablar de las TF de ROS.
	
\end{frame}

\begin{frame}
	\frametitle{Rosbags}
	topic node hz echo Bw etc...
	
\end{frame}

\begin{frame}
	\frametitle{RViz}
	topic node hz echo Bw etc...
	
\end{frame}

\begin{frame}
	\frametitle{Crear mensajes customizados}
	Hablar de las TF de ROS.
	
\end{frame}

\begin{frame}
	\frametitle{Herramientas por línea de comandos útiles}
    topic node hz echo Bw etc...
    ros2 <hit\_tab>
    ros2 topic list
    ros2 bag 
	
\end{frame}


\end{document}
